% ===================== 阅读 Reading =====================
\ptesection{阅读}{Reading · FIB + MCM + RO + MCS}{42}

\vspace{1mm}
{\small\color{gray} 本部分依据「阅读1 / 阅读2」逐题誊录。题型：\textbf{FIB} 读写填空（下拉选词）、\textbf{FIB(R)} 读填空（拖词入空）、
\textbf{MCM} 多选、\textbf{RO} 段落重排、\textbf{MCS} 单选。每题保留：\textbf{文章（含填空/选择答案）+ 中文翻译 + 选项 + 平台中文解析（根据原文 + 解题思路）+ 评分}。}
\vspace{1mm}

\begin{notebox}
{\footnotesize 标注：填空/选择处 {\color{cFix}\ding{51}\,绿=选对}；{\color{cWrong}\ding{55}\,红删除线=选错}，其后 {\footnotesize\color{cFix}(答案: 正确词)} 给正确答案。
选项里**加粗**的是正确项。共 16 题，本卷收录可见的 15 题（缺 Q10，截图未含，见文末）。}
\end{notebox}

% =================== Q1 ===================
\begin{qblock}{Q1\quad FIB \#74\ \ Good Looks in Votes}{读写填空 · 评分 \textcolor{cAccent}{\bfseries 2 / 4}}
\ptesub{文章 Passage（含填空作答）}
\begin{promptbox}
It is tempting to try to prove that good looks win votes, and many academics have tried. The \rok{difficulty} is that beauty is in the eye of the \rno{audience}{beholder}, and you cannot behold a politician's face without a veil of extraneous prejudice getting in the way. Does George Bush possess a disarming grin, or a facetious \rno{binge}{smirk}? It's hard to find anyone who can look at the president without assessing him politically as well as \rok{physically}.
\end{promptbox}
\zh{想证明“颜值能赢得选票”很有诱惑力，许多学者也尝试过。难点在于“情人眼里出西施”——你看一位政客的脸时，总有一层无关的偏见挡在中间。乔治·布什是有一种令人放松的笑容，还是一种假笑？很难找到谁能在看总统时，不在政治层面、同时也在外貌层面对他做出评判。}
\ptesub{选项 Choices}
{\small
1.\ principle,\ idea,\ \textbf{difficulty},\ concept\\
2.\ people,\ \textbf{beholder},\ builder,\ audience\\
3.\ smell,\ complexion,\ \textbf{smirk},\ binge\\
4.\ culturally,\ \textbf{physically},\ economically,\ individually}
\ptesub{解析 Explanation}
\begin{notebox}\footnotesize
\textbf{1.\ difficulty} —— 根据原文：The \_\_\_ is that beauty is in the eye of the beholder…；此处需名词表“难处”，故 difficulty；principle 原则、idea 想法、concept 概念语意不合。\textbf{思路：单句理解}。\\[2pt]
\textbf{2.\ beholder} —— 固定搭配 \emph{beauty is in the eye of the beholder}（情人眼里出西施）；people 人们、builder 建造者、audience 观众均不合。\textbf{思路：固定搭配}。\\[2pt]
\textbf{3.\ smirk} —— 根据原文：Does George Bush possess a disarming grin, or a facetious \_\_\_? 与 grin 并列、需含贬义的名词，故 smirk（假笑）；smell 气味、complexion 肤色、binge 狂欢不合。\textbf{思路：单句理解}。\\[2pt]
\textbf{4.\ physically} —— 根据原文：…assessing him politically as well as \_\_\_；与 politically 并列，按逻辑选 physically；culturally 文化上、economically 经济上、individually 个人地不合。\textbf{思路：上下文逻辑}。
\end{notebox}
\end{qblock}

% =================== Q2 ===================
\begin{qblock}{Q2\quad FIB \#50\ \ Underground Houses}{读写填空 · 评分 \textcolor{cAccent}{\bfseries 2 / 5}}
\ptesub{文章 Passage（含填空作答）}
\begin{promptbox}
Underground houses have many advantages over conventional housing. Unlike conventional homes, they can be built on \rno{overhead}{steep} surfaces and can maximize space in small areas by going below the surface. In addition, the materials excavated in construction can be used in the building process. Underground houses have less surface area so fewer building materials are used, and \rok{maintenance} costs are lower. They are also wind, fire, and earthquake resistant, providing a secure and safe environment in extreme weather. One of the greatest benefits of underground living is energy \rno{consistency}{efficiency}. The earth's subsurface temperature remains stable, so underground dwellings benefit from geothermal mass and heat exchange, staying cool in the summer and warm in the winter. This saves around 80\% in energy costs. By \rno{inventing}{incorporating} solar design this energy bill \rok{can be reduced} to zero, providing hot water and heat to the home all year round.
\end{promptbox}
\zh{地下房屋相比传统住宅有许多优势。与传统住宅不同，它们可以建在陡峭的地面上，并通过向地下发展、在狭小空间里最大化利用面积。此外，施工中挖出的材料还能用于建造过程。地下房屋表面积更小，用的建筑材料更少，维护成本也更低。它们还能抗风、抗火、抗震，在极端天气下提供安全的环境。地下居住最大的好处之一是能源效率：地表以下温度稳定，地下住宅因而受益于地热与热交换，夏凉冬暖，可节省约 80\% 的能源开支。通过引入太阳能设计，这笔能源账单可被降至零，全年为家庭提供热水和供暖。}
\ptesub{选项 Choices}
{\small
1.\ geometric,\ flat,\ overhead,\ \textbf{steep}\\
2.\ \textbf{maintenance},\ heating,\ buoyancy,\ facility\\
3.\ ratio,\ consistency,\ \textbf{efficiency},\ renewal\\
4.\ intriguing,\ initiating,\ \textbf{incorporating},\ inventing\\
5.\ has reduced,\ \textbf{can be reduced},\ can reduce,\ has been reduced}
\ptesub{解析 Explanation}
\begin{notebox}\footnotesize
\textbf{1.\ steep} —— 根据原文：they can be built on \_\_\_ surfaces…；可建在“陡峭”的地面上，故 steep；geometric 几何的、flat 平的、overhead 头顶上的 不合。\textbf{思路：单句理解}。\\[2pt]
\textbf{2.\ maintenance} —— 根据原文：…fewer building materials are used, and \_\_\_ costs are lower；面积小、维修费用低，故 maintenance（维护）。\textbf{思路：单句理解}。\\[2pt]
\textbf{3.\ efficiency} —— 结合下文 \emph{This saves around 80\% in energy costs}（省约 80\% 能耗），故能源“效率”efficiency；ratio 比率、consistency 浓度、renewal 更新 不合。\textbf{思路：上下文理解}。\\[2pt]
\textbf{4.\ incorporating} —— 根据原文：By \_\_\_ solar design… providing hot water…；需“引入/纳入”太阳能设计，故 incorporating；intriguing 引发、initiating 发起、inventing 发明 不合。\textbf{思路：动宾搭配}。\\[2pt]
\textbf{5.\ can be reduced} —— 主语 energy bill 与“减少”是被动关系，且为假设、非已完成，故被动 can be reduced；has reduced、can reduce、has been reduced 不合。\textbf{思路：时态 + 被动语态}。
\end{notebox}
\end{qblock}

% =================== Q3 ===================
\begin{qblock}{Q3\quad FIB \#3\ \ Intelligence Comparison}{读写填空 · 评分 \textcolor{cAccent}{\bfseries 2 / 5}}
\ptesub{文章 Passage（含填空作答）}
\begin{promptbox}
Psychologists have a technique for looking at intelligence that \rok{does} not require the cooperation of the animal involved \dots\ of intelligence. Comparing \rno{with}{across} species is not as simple as generally expected. An elephant will have a larger brain than a human has simply because it is a large beast. \rno{Otherwise}{Instead}, we use the Cephalization index, which compares the size of an animal's brain with the size of its body. Based on the Cephalization index, the brightest animals on the planet are humans, \rok{followed} by great apes, porpoises and elephants. As a general \rno{principle}{rule}, animals that hunt for a living (like canines) are smarter than strict vegetarians (you don't need much intelligence to outsmart a leaf of lettuce). Animals that live in social groups are always smarter and have larger EQ's than solitary animals.
\end{promptbox}
{\footnotesize\color{gray}（注：首句中段在原图分页处未截全。）}
\zh{心理学家有一套观察智力的方法，它不需要相关动物的配合……（与智力有关）。跨物种比较并不像一般以为的那么简单：大象的脑比人大，只是因为它体型大。于是我们改用“脑化指数”，把动物大脑的大小与其身体大小相比较。按脑化指数，地球上最聪明的动物是人类，其次是类人猿、海豚和大象。一般来说，以狩猎为生的动物（如犬科）比纯素食动物更聪明（要智胜一片生菜并不需要多少智力）。群居动物总是比独居动物更聪明、EQ 更高。}
\ptesub{选项 Choices}
{\small
1.\ can,\ do,\ did,\ \textbf{does}\\
2.\ \textbf{across},\ to,\ through,\ with\\
3.\ Then,\ \textbf{Instead},\ Because,\ Otherwise\\
4.\ \textbf{followed},\ follows,\ follow,\ following\\
5.\ theory,\ principal,\ \textbf{rule},\ principle}
\ptesub{解析 Explanation}
\begin{notebox}\footnotesize
\textbf{1.\ does} —— 根据原文：…a technique…that \_\_\_ not require the cooperation of the animal involved；需助动词、第三人称单数现在时，故 does；can/did/do 不匹配。\textbf{思路：主谓一致 + 时态}。\\[2pt]
\textbf{2.\ across} —— 根据原文：Comparing \_\_\_ species…；表“跨/不同物种间”比较，故 across；to/through/with 不合。\textbf{思路：单句理解}。\\[2pt]
\textbf{3.\ Instead} —— 根据原文：\_\_\_, we use the Cephalization index…；与前句形成转折，故副词 Instead；Then/Because/Otherwise 不合。\textbf{思路：单句理解（转折）}。\\[2pt]
\textbf{4.\ followed} —— 根据原文：…humans, \_\_\_ by great apes…；表“其次/紧随其后”，过去分词作后置定语，故 followed；follows/follow/following 不合。\textbf{思路：过去分词}。\\[2pt]
\textbf{5.\ rule} —— 固定搭配 \emph{as a general rule}（一般来说）；theory 理论、principal 校长、principle 法则 不合。\textbf{思路：固定搭配}。
\end{notebox}
\end{qblock}

% =================== Q4 ===================
\begin{qblock}{Q4\quad FIB \#70\ \ Mini Helicopter}{读写填空 · 评分 \textcolor{cAccent}{\bfseries 2 / 5}}
\ptesub{文章 Passage（含填空作答）}
\begin{promptbox}
A mini helicopter modelled on flying tree seeds could soon be flying overhead. Evan Ulrich and colleagues at the University of Maryland in College Park \rno{turned in}{turned to} the biological world for inspiration to build a scaled-down helicopter that could mimic the properties of full-size aircraft. The complex \rok{design} of full-size helicopters gets less efficient when shrunk, meaning that standard mini helicopters expend most of their power simply fighting to stay stable in the air. The researchers realized that a simpler aircraft designed to stay stable passively would use much less power and reduce manufacturing costs to boot. It turns out that nature \rno{is beating}{had beaten} them to it. The seeds of trees have a single-blade structure that \rok{allows} them to fly far away and drift safely to the ground. These seeds need no engine to \rno{fly}{spin} through the air, thanks to a process called autorotation. By analyzing the behavior of the samara with high-speed cameras, Ulrich and his team were able to copy its design.
\end{promptbox}
\zh{一种以飞行树种子为原型的微型直升机可能很快就会在头顶飞行。马里兰大学帕克分校的 Evan Ulrich 及同事转向生物界寻找灵感，想造一架能模仿全尺寸飞行器特性的缩小版直升机。全尺寸直升机复杂的设计在缩小后效率会下降，这意味着普通微型直升机的大部分动力都耗在勉力维持空中稳定上。研究者意识到，一架靠被动方式保持稳定的更简单飞行器会省电得多，还能降低制造成本。事实证明，大自然早已抢先一步：树的种子有一种单叶片结构，能让它们飞得很远并安全飘落到地面。借助一种叫“自转”的过程，这些种子无需发动机就能在空中旋转。通过用高速摄像机分析翅果（samara）的运动，Ulrich 团队得以复制其设计。}
\ptesub{选项 Choices}
{\small
1.\ \textbf{turned to},\ turned for,\ turned in,\ turned off\\
2.\ overhaul,\ gauge,\ imagination,\ \textbf{design}\\
3.\ is beating,\ was beaten,\ \textbf{had beaten},\ beaten\\
4.\ had allowed,\ allowed,\ \textbf{allows},\ will allow\\
5.\ \textbf{spin},\ fluctuate,\ drift,\ fly}
\ptesub{解析 Explanation}
\begin{notebox}\footnotesize
\textbf{1.\ turned to} —— 根据原文：…colleagues \_\_\_ the biological world for inspiration…；转向自然界寻求灵感，故 turned to（求助于）；turned in/for/off 不合。\textbf{思路：单句理解}。\\[2pt]
\textbf{2.\ design} —— 根据原文：The \_\_\_ of full-size helicopters gets less efficient when shrunk…；指“设计”，故 design；overhaul 检修、gauge 仪表、imagination 想象 不合。\textbf{思路：上下文理解}。\\[2pt]
\textbf{3.\ had beaten} —— 根据原文：It turns out that nature \_\_\_ them to it；\emph{beat sb. to it}（抢先一步），与 realized 等过去时呼应，表“过去的过去”，故过去完成 had beaten；is beating/was beaten/beaten 不合。\textbf{思路：时态}。\\[2pt]
\textbf{4.\ allows} —— 根据原文：…a single-blade structure that \_\_\_ them to fly far away…；定语从句谓语、一般现在时，故 allows；had allowed/allowed/will allow 时态不合。\textbf{思路：时态}。\\[2pt]
\textbf{5.\ spin} —— 根据原文：These seeds need no engine to \_\_\_ through the air, thanks to a process called autorotation；autorotation（自旋），故 spin（旋转）；fluctuate 波动、drift 漂流、fly 飞 不合。\textbf{思路：单句理解}。
\end{notebox}
\end{qblock}

% =================== Q5 ===================
\begin{qblock}{Q5\quad FIB \#290\ \ Power Mix}{读写填空 · 评分 \textcolor{cAccent}{\bfseries 3 / 5}}
\ptesub{文章 Passage（含填空作答）}
\begin{promptbox}
Imagine a time in the not too distant future when your power comes from a seamless mix of renewable energy and traditional sources. It is delivered by a grid that manages thousands of windmills and hundreds of thousands of customers. Computer \rno{controlling}{controlled}, the grid is able to manage instant variations in supply and demand and provides a real time power balance. Far more complex than anything \rok{in} existence today, it is the smart grid. This technology is a new frontier in power supply and seen as a green solution to current outdated management systems. When introduced, smart grids will result in energy savings and will allow consumers a choice in their electricity charges and to be able to select the cheapest time \rno{pins}{slots}. The difficulty for the energy industry is that smart grids \rok{do not exist} in reality and the power companies cannot experiment with existing supplies. Without an actual grid to conduct research on, Professor Wu has had to design a simulated laboratory including input from theoretical wind generators and solar panels to feed into a constantly operating system. For an authentic approach researchers built various types of equipment failures \rok{into} the grid to test the system. And it works.
\end{promptbox}
\zh{设想在不远的将来，你的电力来自可再生能源与传统能源的无缝组合，由一张管理着成千上万台风车、数十万用户的电网输送。在计算机控制下，电网能管理供需的瞬时变化，提供实时的电力平衡。它远比当今现存的任何系统都复杂——这就是智能电网。这项技术是电力供应的新前沿，被视为解决当前陈旧管理系统的绿色方案。智能电网一旦引入，将带来节能，并让消费者在电费上有选择权、能挑选最便宜的时段。能源行业的难处在于：智能电网在现实中并不存在，电力公司无法用现有供电做实验。由于没有真实电网可供研究，Wu 教授只得设计一个模拟实验室，引入理论风力发电机和太阳能板的输入，接入一个持续运行的系统。为贴近真实，研究者在电网中植入各种类型的设备故障来测试系统。结果——它奏效了。}
\ptesub{选项 Choices}
{\small
1.\ \textbf{controlled}\ {\footnotesize\color{gray}（其余选项在原图分页处未截全）}\\
2.\ \textbf{in}\ {\footnotesize\color{gray}（同上）}\\
3.\ cuts,\ pins,\ points,\ \textbf{slots}\\
4.\ does not exist,\ \textbf{do not exist},\ are not existing,\ not exist\\
5.\ \textbf{into},\ of,\ onto,\ above}
\ptesub{解析 Explanation}
\begin{notebox}\footnotesize
\textbf{1.\ controlled} —— 根据原文：Computer \_\_\_, the grid is able to…（独立主格结构）；Computer 与 control 为被动关系，用过去分词 controlled。\textbf{思路：独立主格 / 非谓语}。\\[2pt]
\textbf{2.\ in} —— 固定搭配 \emph{in existence}（现存的）：anything \_\_\_ existence today，故 in。\textbf{思路：介词 / 固定搭配}。\\[2pt]
\textbf{3.\ slots} —— 根据原文：…select the cheapest time \_\_\_；\emph{time slots}（时段），故 slots；cuts/pins/points 不合。\textbf{思路：词组搭配}。\\[2pt]
\textbf{4.\ do not exist} —— 根据原文：…smart grids \_\_\_ in reality…；主语复数、一般现在时否定，故 do not exist；does not exist/are not existing/not exist 不合。\textbf{思路：主谓一致}。\\[2pt]
\textbf{5.\ into} —— 根据原文：…built various types of equipment failures \_\_\_ the grid…；表“放入…其中”，故 into（与上文 feed into 呼应）；of/onto/above 不合。\textbf{思路：介词}。
\end{notebox}
\end{qblock}

% =================== Q6 ===================
% （待填充）
